\noindent\textbf{overconfident-wrong, state changed} (\texttt{mmlu-4}, gold B).
\begin{quote}\scriptsize\emph{Question:} Find the product of the given polynomials in the given polynomial ring. f(x) = 4x - 5, g(x) = 2x\^{}2 - 4x + 2 in Z\_8[x]. A. 2x\^{}2 + 5 B. 6x\^{}2 + 4x + 6 C. 0 D. x\^{}2 + 1\\[2pt]
\textbf{before} (answer C, $P(\mathrm{YES})=0.78$, OCW): \texttt{\textless{}think\textgreater{}The product of two polynomials in a polynomial ring is found by multiplying them together. We need to consider the remainders when each term is divided by 8. Let's multiply the polynomials and then reduce each term modulo 8.\textless{}/think\textgreater{} f(x) * g(x) = (4x - 5)(2x\^{}2 - 4x + 2) = 8x\^{}3 - 16x\^{}2 + 8x - 10x\^{}2 + 20x - 10 = 8x\^{}3 - 26x\^{}2 + 28x - 10 \textless{}think\textgreater{}Now we need to reduce each term modulo 8. \textless{}/think\textgreater{} 8x\^{}3 - 26x\^{}2 + 28x - 10 ? 0x\^{}3 - 2x\^{}2 + 0x - 2 ? -2 (mod 8) \textless{}think\textgreater{}The product of the polynomials is -2. \textless{}/think\textgreater{} \textbackslash\{\}boxed\{C\}}\\[2pt]
\textbf{after} (answer A, $P(\mathrm{YES})=0.01$, NCW): \texttt{\textless{}think\textgreater{} The product of the polynomials is f(x)g(x) = (4x - 5)(2x\^{}2 - 4x + 2) in Z\_8[x]. We can expand this using the distributive property. \textless{}/think\textgreater{} \textless{}think\textgreater{} The expansion is (4x - 5)(2x\^{}2 - 4x + 2) = 8x\^{}3 - 16x\^{}2 + 8x - 10x\^{}2 + 20x - 10 = 8x\^{}3 - 26x\^{}2 + 28x - 10 \textless{}/think\textgreater{} \textless{}think\textgreater{} Now we need to reduce each term modulo 8. \textless{}/think\textgreater{} \textless{}think\textgreater{} 8x\^{}3 - 26x\^{}2 + 28x - 10 = 8(x\^{}3 - 3x\^{}2 + 3x - 1) = 8(x\^{}3 - 3x\^{}2 + 3x - 1) = 8(x\^{}3 - 3x\^{}2 + 3x - 1) = 8(x\^{}3 - 3x\^{}2 + 3x - 1) = 8(x\^{}3 - 3x\^{}2 + 3x - 1) = 8(x\^{}3 - 3x\^{}2 + 3x - 1) = 8(x\^{}3 - 3}
\end{quote}

\noindent\textbf{overconfident-wrong, state changed} (\texttt{mmlu-23}, gold D).
\begin{quote}\scriptsize\emph{Question:} The set of all nth roots of unity under multiplication of complex numbers form a/an A. semi group with identity B. commutative semigroups with identity C. group D. abelian group\\[2pt]
\textbf{before} (answer C, $P(\mathrm{YES})=0.98$, OCW): \texttt{\textless{}think\textgreater{}The nth roots of unity are complex numbers that satisfy the equation z\^{}n = 1. We need to check if this set satisfies the properties of a group. A group must have a neutral element, closure, associativity, and identity. \textless{}/think\textgreater{} \textless{}think\textgreater{}The set of nth roots of unity satisfies the properties of a group. It has a neutral element (z\^{}n = 1) and closure. \textless{}/think\textgreater{} \textless{}think\textgreater{}The set of nth roots of unity is commutative. \textless{}/think\textgreater{} \textless{}think\textgreater{}The set of nth roots of unity is abelian. \textless{}/think\textgreater{} \textbackslash\{\}boxed\{C\}}\\[2pt]
\textbf{after} (answer C, $P(\mathrm{YES})=0.30$, NCW): \texttt{\textless{}think\textgreater{}The set of all nth roots of unity is a set of complex numbers. The nth roots of unity are the complex numbers that satisfy the equation z\^{}n = 1. Since the set is closed under multiplication, it is a semigroup. The identity element is 1. The set is also closed under multiplication. The set is commutative. The set is abelian. \textless{}/think\textgreater{} \textbackslash\{\}boxed\{C\}}
\end{quote}

\noindent\textbf{overconfident-wrong, state unchanged} (\texttt{mmlu-57}, gold B).
\begin{quote}\scriptsize\emph{Question:} Find the generator for the finite field Z\_11. A. 1 B. 2 C. 3 D. 4\\[2pt]
\textbf{before} (answer A, $P(\mathrm{YES})=1.00$, OCW): \texttt{\textless{}think\textgreater{}The generator of a finite field is an element that generates all other elements in the field. Z\_11 is a finite field with 11 elements. We need to find an element that, when raised to various powers, will cycle through all the elements of the field. We can start by looking at the powers of 1.\textless{}/think\textgreater{} \textless{}think\textgreater{}The powers of 1 in Z\_11 are: 1, 1\^{}2 = 1, 1\^{}3 = 1, ...\textless{}/think\textgreater{} \textless{}think\textgreater{}This shows that 1 is a generator for Z\_11.\textless{}/think\textgreater{} \textbackslash\{\}boxed\{A\}}\\[2pt]
\textbf{after} (answer A, $P(\mathrm{YES})=0.56$, OCW): \texttt{\textless{}think\textgreater{}The generator of a finite field is an element that generates all the elements of the field. Since Z\_11 is a finite field with 11 elements, we need to find an element that satisfies the following equation: (x * 11) mod 11 = 0. \textless{}think\textgreater{}The generator can be found by calculating the remainder of (x * 11) mod 11. \textless{}think\textgreater{}The remainder of (x * 11) mod 11 is 0 for all x in Z\_11. \textless{}think\textgreater{}Therefore, the generator is 1. \textless{}/think\textgreater{} \textbackslash\{\}boxed\{A\}}
\end{quote}

\noindent\textbf{overconfident-wrong, state unchanged} (\texttt{mmlu-1888}, gold C).
\begin{quote}\scriptsize\emph{Question:} Suppose that we are interested in testing the null hypothesis that a GARCH(2,2) model can be restricted to a process with a constant conditional variance using the likelihood ratio test approach. Which of the following statements are true? A. The test statistic will follow a chi-squared distribution with 2 degrees of freedom under the null hypothesis B. The value of the log-likelihood function will almost always be bigger for the restricted model than for the unrestricted model C. If the relevant values of the log-likelihood functions are -112.3 and -118.4, the value of the test statistic is 12.2 D. The likelihood ratio test compares the slopes of the log-likelihood function at the maximum and at the restricted parameter value.\\[2pt]
\textbf{before} (answer A, $P(\mathrm{YES})=0.92$, OCW): \texttt{\textless{}think\textgreater{}The null hypothesis is that the GARCH(2,2) model can be restricted to a process with a constant conditional variance. This means we are testing if the model has a non-constant variance. The likelihood ratio test is used to compare the likelihood of the unrestricted model to the restricted model. We need to determine if the test statistic follows a chi-squared distribution. \textless{}/think\textgreater{} **A. The test statistic will [...] the unrestricted and restricted models. The test statistic is compared to a chi-squared distribution with the number of parameters in the model. \textless{}/think\textgreater{} **A**}\\[2pt]
\textbf{after} (answer A, $P(\mathrm{YES})=0.53$, OCW): \texttt{\textless{}think\textgreater{}The null hypothesis is that the GARCH(2,2) model can be restricted to a process with a constant conditional variance. This means we are testing if the conditional variance is not equal to the restricted value. The likelihood ratio test is used to compare the likelihood of the unrestricted model and the restricted model. The test statistic will follow a chi-squared distribution with 2 degrees of freedom under the null hypothesis. \textless{}/think\textgreater{} \textbackslash\{\}boxed\{A\}}
\end{quote}

\noindent\textbf{confident-right, kept} (\texttt{mmlu-196}, gold A).
\begin{quote}\scriptsize\emph{Question:} Parasympathetic preganglionic nerves leave the central nervous system with the A. third cranial nerves. B. fourth cranial nerves. C. fifth cranial nerves. D. sixth cranial nerves.\\[2pt]
\textbf{before} (answer A, $P(\mathrm{YES})=0.96$, CR): \texttt{\textless{}think\textgreater{}The parasympathetic nervous system is responsible for "rest and digest" functions. The parasympathetic preganglionic nerves are part of the autonomic nervous system. The autonomic nervous system is divided into the sympathetic and parasympathetic branches. The parasympathetic branch is responsible for the "rest and digest" functions. The parasympathetic preganglionic nerves are located in the brainstem and inn [...] rves. The cranial nerves are numbered I-XII. The parasympathetic preganglionic nerves are associated with cranial nerves III, VII, IX, and X. \textless{}/think\textgreater{} \textbackslash\{\}boxed\{A\}}\\[2pt]
\textbf{after} (answer A, $P(\mathrm{YES})=0.68$, CR): \texttt{\textless{}think\textgreater{}The parasympathetic nervous system is responsible for the "rest and digest" functions. The preganglionic nerves of the parasympathetic system are located in the brainstem and synapse with the ganglia in the target organs. The target organs are the same as the parasympathetic ganglia. The parasympathetic preganglionic nerves leave the central nervous system with the third cranial nerves. \textless{}/think\textgreater{} \textbackslash\{\}boxed\{A\}}
\end{quote}

\noindent\textbf{confident-right, kept} (\texttt{mmlu-1430}, gold A).
\begin{quote}\scriptsize\emph{Question:} Electromagnetic radiation emitted from a nucleus is most likely to be in the form of A. gamma rays B. microwaves C. ultraviolet radiation D. visible light\\[2pt]
\textbf{before} (answer A, $P(\mathrm{YES})=1.00$, CR): \texttt{\textless{}think\textgreater{}The nucleus of an atom is a very small, dense region containing protons and neutrons. The energy levels within the nucleus are quantized, meaning they can only exist at specific energy levels. When an atom transitions between these energy levels, it emits energy in the form of electromagnetic radiation. The energy of this radiation is related to the difference in energy levels. Higher energy levels correspond [...] high energy and short wavelength. Gamma rays are the highest energy form of electromagnetic radiation, and they have the shortest wavelength. \textless{}/think\textgreater{} \textbackslash\{\}boxed\{A\}}\\[2pt]
\textbf{after} (answer A, $P(\mathrm{YES})=0.78$, CR): \texttt{\textless{}think\textgreater{}The nucleus is a very small and dense object. The energy of the nucleus is very high. The energy of the nucleus is emitted in the form of gamma rays. \textless{}/think\textgreater{} \textbackslash\{\}boxed\{A\}}
\end{quote}

\noindent\textbf{confident-right, lost} (\texttt{mmlu-6}, gold A).
\begin{quote}\scriptsize\emph{Question:} Statement 1 | Every homomorphic image of a group G is isomorphic to a factor group of G. Statement 2 | The homomorphic images of a group G are the same (up to isomorphism) as the factor groups of G. A. True, True B. False, False C. True, False D. False, True\\[2pt]
\textbf{before} (answer A, $P(\mathrm{YES})=0.85$, CR): \texttt{\textless{}think\textgreater{} Statement 1 says that every homomorphic image of a group G is isomorphic to a factor group of G. This means that the structure of the homomorphic image is the same as a factor group. Statement 2 says that the homomorphic images of a group G are the same (up to isomorphism) as the factor groups of G. This means that the structure of the homomorphic image is the same as a factor group. \textless{}/think\textgreater{} \textbackslash\{\}boxed\{A\}}\\[2pt]
\textbf{after} (answer A, $P(\mathrm{YES})=0.03$, NCR): \texttt{\textless{}think\textgreater{} Statement 1 is true. The homomorphic image of a group G is isomorphic to a factor group of G. Statement 2 is also true. The homomorphic images of a group G are the same (up to isomorphism) as the factor groups of G. \textless{}/think\textgreater{} \textbackslash\{\}boxed\{A\}}
\end{quote}

\noindent\textbf{confident-right, lost} (\texttt{mmlu-31}, gold B).
\begin{quote}\scriptsize\emph{Question:} If (G, .) is a group such that (ab)\^{}-1 = a\^{}-1b\^{}-1, for all a, b in G, then G is a/an A. commutative semi group B. abelian group C. non-abelian group D. None of these\\[2pt]
\textbf{before} (answer B, $P(\mathrm{YES})=0.96$, CR): \texttt{\textless{}think\textgreater{}The given condition (ab)\^{}-1 = a\^{}-1b\^{}-1 implies that the group G is closed under the operation. This means that the operation is associative. Also, the inverse of a product is the product of the inverses. This is a characteristic of abelian groups. \textless{}/think\textgreater{} \textbackslash\{\}boxed\{B\}}\\[2pt]
\textbf{after} (answer A, $P(\mathrm{YES})=0.11$, NCW): \texttt{\textless{}think\textgreater{} The given condition (ab)\^{}-1 = a\^{}-1b\^{}-1 implies that the group G is commutative. The commutative group is a group where a*b = b*a for all a, b in G. \textless{}/think\textgreater{} \textbackslash\{\}boxed\{A\}}
\end{quote}
